%----------------------------------------------------------------------------------------
% METADATA BLOCK (YAML)  — see inbox/README.md for every field
%----------------------------------------------------------------------------------------
% uid: ""
% school: "lycee"
% level: "2eme"
% section: ["Sciences", "Informatique"]
% chapter: "Suites arithmétiques et géométriques"
% title: "Suite arithmétique et somme des premiers termes"
% difficulty: "Facile"
% points: 6
% tags: ["suites-arithmetiques", "terme-general", "somme"]
% provenance: "own"
% dressing: "deux termes donnés"
% rung: 2
% batch: "2026-11-18-suites-2sc-A"
%----------------------------------------------------------------------------------------

% Everything above this line is metadata. Below: plain LaTeX content only
% (no tcolorbox, no custom environments, no handout styling).
% Macros defined here are expanded by the converter; keep them simple.
\newcommand{\N}{\mathbb{N}}

% ---STATEMENT---
Soit $(u_n)$ une suite arithmétique définie sur $\N$ telle que $u_2 = 7$ et $u_5 = 13$.

\begin{enumerate}
  \item
  \begin{enumerate}
    \item Déterminer la raison $r$ de la suite $(u_n)$.
    \item Calculer $u_0$ puis exprimer $u_n$ en fonction de $n$.
  \end{enumerate}
  \item Pour tout $n \in \N$, on pose $S_n = u_0 + u_1 + \dots + u_n$.
  \begin{enumerate}
    \item Montrer que $S_n = (n+1)(n+3)$.
    \item Déterminer l'entier $n$ tel que $S_n = 195$.
  \end{enumerate}
\end{enumerate}

% ---SOLUTION---
\begin{enumerate}
  \item
  \begin{enumerate}
    \item Dans une suite arithmétique, $u_p = u_q + (p-q)\,r$.

    Avec $p = 5$ et $q = 2$ : $u_5 = u_2 + 3r$, donc $13 = 7 + 3r$.

    D'où $3r = 6$ et donc $r = 2$.
    \item On utilise $u_n = u_0 + n r$ avec $n = 2$ : $7 = u_0 + 4$, donc $u_0 = 3$.

    Ainsi $u_n = 3 + 2n$ pour tout $n \in \N$.
  \end{enumerate}
  \item
  \begin{enumerate}
    \item La somme de $n+1$ termes consécutifs d'une suite arithmétique vaut
    $$S_n = (n+1)\,\frac{u_0 + u_n}{2}.$$

    Ici $u_0 + u_n = 3 + (3 + 2n) = 2n + 6$, donc
    $$S_n = (n+1)\,\frac{2n+6}{2} = (n+1)(n+3).$$
    \item On résout $(n+1)(n+3) = 195$, c'est-à-dire $n^2 + 4n - 192 = 0$.

    Le discriminant est $\Delta = 16 + 768 = 784 = 28^2$, donc $n = \dfrac{-4 + 28}{2} = 12$ ou $n = \dfrac{-4-28}{2} = -16$.

    Comme $n \in \N$, on garde $n = 12$. Vérification : $13 \times 15 = 195$.
  \end{enumerate}
\end{enumerate}

% alt: Les points (n, u_n) pour n de 0 à 5 sont alignés sur une droite de pente 2.
\begin{center}
\begin{tikzpicture}[x=1cm,y=0.3cm]
  \draw[->] (-0.3,0) -- (6,0) node[right] {$n$};
  \draw[->] (0,-0.3) -- (0,14) node[above] {$u_n$};
  \foreach \i in {0,...,5} {
    \fill (\i,{3+2*\i}) circle (2pt);
    \draw[dotted] (\i,0) -- (\i,{3+2*\i});
    \node[below] at (\i,0) {\small \i};
  }
  \draw (0,3) -- (5,13);
\end{tikzpicture}
\end{center}
